\section{What Makes Questions Prone to Overthinking?}

\subsection{Feature Analysis}

We extract features from each question and compare overthinking vs stable samples.

\paragraph{Question length.} Overthinking samples have significantly shorter questions:
\begin{itemize}[nosep]
\item Overthinking: 39.2 words (avg)
\item Stable correct: 45.1 words
\item Improved (256✗→512✓): 48.8 words
\end{itemize}

This suggests \textbf{simpler questions are more prone to overthinking}.

\paragraph{Numerical complexity.} We count numbers in each question:
\begin{itemize}[nosep]
\item Overthinking: 6.3 numbers (avg)
\item Stable: 7.8 numbers
\item Improved: 8.9 numbers
\end{itemize}

\subsection{Predictive Model}

We train a logistic regression classifier to predict overthinking risk from question features:
$$P(\text{overthink} | x) = \sigma(w^\top \phi(x))$$

where $\phi(x)$ includes: question length, number count, max number value, presence of fractions/percentages.

\textbf{Results:} The classifier achieves 68\% accuracy, significantly above chance (50\%).
Feature importance analysis reveals question length is the strongest predictor (coefficient: -0.42).

\subsection{Implications}

The predictability of overthinking suggests it is a \emph{systematic} phenomenon, not random noise.
Short, simple questions are at highest risk---precisely the cases where extended reasoning seems least necessary.
